% 分野別類題: 一次同次線形偏微分方程式（P u_x + Q u_y = 0）
% 特性曲線法で解く基本形 3 問。
% コンパイル: TEXINPUTS="../../../00_common:$TEXINPUTS" latexmk -xelatex problems.tex
\documentclass[a4paper,11pt]{article}
\usepackage{preamble}

\begin{document}

\kamokumei{類題演習 --- 一次同次線形偏微分方程式}

\shijibun{次の偏微分方程式の一般解を求めよ。導出過程を示せ。（ヒント: $P(x,y)\,\dfrac{\partial u}{\partial x} + Q(x,y)\,\dfrac{\partial u}{\partial y} = 0$ の一般解は、特性方程式 $\dfrac{dx}{P} = \dfrac{dy}{Q}$ の解 $\varphi(x,y) = C$ を用いて $u = F(\varphi(x,y))$（$F$ は任意関数）と表せる。）}

\shomon{問1.}{次の偏微分方程式の一般解を求めよ。}
\[
x\frac{\partial u}{\partial x} - y\frac{\partial u}{\partial y} = 0
\]

\shomon{問2.}{次の偏微分方程式の一般解を求めよ。}
\[
y\frac{\partial u}{\partial x} + x\frac{\partial u}{\partial y} = 0
\]

\shomon{問3.}{次の偏微分方程式の一般解を求め、さらに条件 $u(0, y) = \sin y$ を満たす解を求めよ。}
\[
\frac{\partial u}{\partial x} + 2x\frac{\partial u}{\partial y} = 0
\]

\end{document}
